\documentclass[10pt,a4paper]{amsart}
\usepackage[margin=19mm]{geometry}
\usepackage{amsmath,amssymb,mathtools}
\usepackage[scr=rsfs,cal=euler]{mathalfa}
\usepackage{xfrac}
\usepackage[unicode,hidelinks,bookmarks=false]{hyperref}
\newcommand{\ie}{i.e.\ }
\newcommand{\cf}{cf.\ }
\newcommand{\resp}{respectively\ }
\newcommand{\Subsection}{\subsection}
\providecommand{\nobreakdash}{\nobreak\hbox{-}}
\setlength{\emergencystretch}{2em}
\multlinegap0pt
\usepackage{amssymb}
\usepackage{xcolor}
\usepackage{tikz}
\usepackage{caption}
\usepackage[shortlabels]{enumitem}
\usepackage{cancel}
\newtheorem{assumption}{Assumption}
\newtheorem{theorem}{Theorem}
\newtheorem{lemma}{Lemma}
\usepackage{cleveref}
\crefname{assumption}{Assumption}{Assumptions}
\crefname{theorem}{Theorem}{Theorems}
\crefname{lemma}{Lemma}{Lemmas}
\DeclareMathOperator{\E}{\mathbb E}
\newcommand{\R}{\mathbb{R}}
\newcommand{\dd}{\ensuremath{\,\mathrm d}}
\title{Functional SDE approximation inspired by a deep operator network architecture}
\author{Martin Eigel, Charles Miranda}
\date{Source: arXiv:2402.03028v2 (CC BY 4.0)}
\begin{document}
\maketitle

\noindent\emph{Source and licence.} Functional SDE approximation inspired by a deep operator network architecture. Version arXiv:2402.03028v2 (CC BY 4.0), distributed by arXiv under the Creative Commons Attribution 4.0 International licence, http://creativecommons.org/licenses/by/4.0/, as recorded in the arXiv licence metadata for this exact version and shown on the abstract page https://arxiv.org/abs/2402.03028v2. Full licence notice: \texttt{licenses/arxiv-2402.03028-CC-BY-4.0.txt}.\par\smallskip

\noindent\textit{Editorial scope.} Counted unit: the article's lemma lemma:ub\_approx\_error (source0.tex lines 756--763). The article's complete original proof of it is delivered with it (source0.tex lines 1144--1176, one \texttt{proof} environment of 33 lines, which the article places away from the statement). No mathematics is written, repaired or re-derived: the statement and the proof are the article's own lines, in the article's own notation, and no retained line is substituted or re-flowed.  Retained uncounted as the source-local context the proof needs: lines 145--148; lines 215--221; lines 284--299; lines 539--541; lines 675--679; lines 742--753.  Nothing was omitted from any retained span, so the package carries no omission record.  No retained line contains a citation, so the wrapper carries no bibliography and no bibliography entry.  No retained line contains a figure environment, an image-inclusion command or drawing code, so no image is delivered and the package declares no asset dependency.  Editorial formatting only, in addition: an AMS \texttt{amsart} preamble that restates the article's own declarations and macros used by the retained lines, namely the environments \texttt{assumption}, \texttt{theorem}, \texttt{lemma} on separate counters, together with the standard packages those lines need.  The retained environments print in this document as Assumption 1, Theorem 1, Lemma 1 in order of appearance, whereas the article prints the same items with the numbering of its own full document.  The complete native source of the article as distributed in the arXiv e-print for this exact version is bundled unchanged as \texttt{sources/154039/source0.tex}.
\begin{equation}
    \dd X_t = \mu(t, X_t)\dd t + \sigma(t, X_t)\dd W_t, \quad\text{with}~X_0=x_0.
    \label{eq:sde}
\end{equation}

\begin{assumption}\label{assump:uniqueness_sde}
    Let $T>0$ and $\mu:[0,T]\times \R^d \to \R^d$, $\sigma : [0,T] \times \R^d \to \R^{d\times m}$ be $\mathcal F$-measurable functions satisfying,
    \begin{enumerate}
        \item linear growth: $\forall x \in \R^d, t \in [0,T], \|\mu(t,x)\|_2 + \|\sigma(t,x)\|_{\text{Fro}} \leq K(1+\|x\|_2)$,
        \item uniform Lipschitz continuity: $\forall x,y \in \R^d, t \in [0,T], \|\mu(t,x)-\mu(t,y)\|_2 + \|\sigma(t,x)-\sigma(t,y)\|_\emph{Fro} \leq K\|x-y\|_2$
    \end{enumerate}
\end{assumption}

\begin{theorem}[{$L^2$-error of the Wiener chaos truncation}]\label{theorem:l2_error_truncation}
    Given \cref{assump:uniqueness_sde} and let $(X_t)_{t \in [0,T]}$ satisfy \cref{eq:sde}.
    Moreover, assume that $\mu,\sigma \in C^{1,\infty}([0,T]\times \R)$ such that
    \begin{equation*}
        \left|\frac{\partial^{\ell+m}\mu}{\partial t^\ell \partial x^m}(t, x) - \frac{\partial^{\ell+m}\mu}{\partial t^\ell \partial x^m}(t, y) \right| + \left|\frac{\partial^{\ell+m}\sigma}{\partial t^\ell \partial x^m}(t, x) - \frac{\partial^{\ell+m}\sigma}{\partial t^\ell \partial x^m}(t, y) \right| \leq K|x-y|,\quad K>0,
    \end{equation*}
    for $t \in [0,T],\ x,y \in \R$, any $\ell \in \{0,1\}$, and $m \geq 0$.
    Then it holds that
    \begin{equation}
        \E[(X_t^{p,k} - X_t)^2] \leq C(1+x_0^2)\left(\frac{1}{(p+1)!} + \sum_{\ell=k+1}^\infty \left(E_\ell(t)^2 + \int_0^t E_\ell^2(\tau)\dd\tau\right) \right), \label{eq:l2_error_truncation}
    \end{equation}
    where $C = C(t, K)$ is a positive constant and the function $E_\ell(t)$ is defined by
    \begin{equation}
        E_\ell(t) := \int_0^t e_\ell(s)\dd s. \label{eq:int_basis}
    \end{equation}
\end{theorem}

\begin{equation}
    X_t(\omega) = \sum_{k \geq 0}\sum_{|\alpha|=k} x_\alpha(t) \Psi_\alpha(\omega), \label{eq:sde_params}
\end{equation}

    \begin{align}
        \hat E_\text{Trunc} &:= \left(\int_0^T \E[|X_t - \sum_{j=1}^p x_{k_j^*}(t)\Psi_{k_j^*}|^2]\dd t\right)^{1/2},\label{eq:trunc_error}\\
        \hat E_\text{Approx} &:= \left(\int_0^T \E[|\sum_{j=1}^p x_{k_j^*}(t)\Psi_{k_j^*} - \sum_{j=1}^p x_{k_j^*}(t)\widetilde{\Psi_j}|^2]\dd t\right)^{1/2},\label{eq:approx_error}\\
        \hat E_\text{Recon} &:= \left(\int_0^T \E[|\sum_{j=1}^p x_{k_j^*}(t)\widetilde{\Psi_j} - \sum_{j=1}^p \widetilde{x_j}(t)\widetilde{\Psi_j}|^2]\dd t\right)^{1/2}.\label{eq:reconstruct_error}
    \end{align}

\begin{lemma}\label{lemma:decay_squared_coef}
    Let \cref{assump:uniqueness_sde} be satisfied, and let $\mu,\sigma$ satisfy the assumptions of \cref{theorem:l2_error_truncation}.
    Consider the Wiener chaos expansion~\cref{eq:sde_params} of $X_t$
    \begin{equation*}
        X_t = \sum_{n=0}^\infty \sum_{|\alpha|=n} x_\alpha(t) \Psi_\alpha.
    \end{equation*}
    Then, for $m \in \mathbb N$ we have
    \begin{equation*}
        \sum_{\ell = 0}^m \sum_{|\alpha|=\ell} x_\alpha(t)^2 \leq (1+x_0^2)\left(e^{Ct e^{Ct}} - \frac{(Cte^{Ct})^{m+1}}{(m+1)!}\right),
    \end{equation*}
    where $t \in [0, T]$, $C = C(K, T)$ is a constant that depends only on $T$ and the regularity of $\mu$ and $\sigma$.
\end{lemma}

\begin{lemma}[{Upper bound of approximation error}]\label{lemma:ub_approx_error}
    Let \cref{assump:uniqueness_sde} be satisfied and let $\mu,\sigma$ satisfy the assumptions of \cref{theorem:l2_error_truncation}.
    Then, the \emph{approximation error}~\cref{eq:approx_error} satisfies
    \begin{align*}
        \hat E_\text{Approx}^2 &\leq \min_{(m,n) \in J_p} \E\left[\sum_{j=1}^p |\Psi_{k_j^*}-\widetilde{\Psi_j}|^2)\right]\\ &\times (1+x_0^2)\left(A(K,T) - \frac{1}{(m+1)!}\frac{T(CT)^{2(m+1)}}{2(m+1)+1}\left(1+\frac{1}{CT}\right)^{m+1}\right),
    \end{align*}
    with constants $A(K,T):=\int_0^T e^{Cte^{Ct}} \dd t$ and $C=C(K,T)$.
\end{lemma}

\begin{proof}[Proof of \cref{lemma:ub_approx_error}]
    By the Cauchy-Schwarz inequality, it follows that
    \begin{align*}
        \left|\sum_{j=1}^p x_{k_j^*}(t)\Psi_{k_j^*} - \sum_{j=1}^p x_{k_j^*}(t)\widetilde{\Psi_j}\right|^2 &= \left|\sum_{j=1}^p x_{k_j^*}(t)(\Psi_{k_j^*}-\widetilde{\Psi_j})\right|^2\\
        &\leq \left(\sum_{j=1}^p x_{k_j^*}(t)^2\right)\left(\sum_{j=1}^p |\Psi_{k_j^*}-\widetilde{\Psi_j}|^2\right).
    \end{align*}
    Taking the expectation and integrating with respect to $t$ leads to
    \begin{align*}
        \int_0^T \E\left[\left|\sum_{j=1}^p x_{k_j^*}(t)\Psi_{k_j^*} - \sum_{j=1}^p x_{k_j^*}(t)\widetilde{\Psi_j}\right|^2\right]\dd t &\leq \int_0^T \left(\sum_{j=1}^p x_{k_j^*}(t)^2\right)\E\left[\sum_{j=1}^p |\Psi_{k_j^*}-\widetilde{\Psi_j}|^2\right]\dd t.
    \end{align*}
    The next step is to find an upper bound of $\int_0^T \left(\sum_{j=1}^p x_{k_j^*}(t)^2\right) \dd t$.
    By \cref{lemma:decay_squared_coef}, with $(m,n) \in J_p$ we have
    \begin{equation*}
        \sum_{\ell = 0}^m \sum_{|\alpha|=\ell} x_\alpha(t)^2 \leq (1+x_0^2)\left(e^{Ct e^{Ct}} - \frac{(Cte^{Ct})^{m+1}}{(m+1)!}\right).
    \end{equation*}
    Integrating with respect to $t$ results in
    \begin{align*}
        \int_0^T \sum_{\ell=0}^m\sum_{|\alpha|=\ell}x_\alpha(t)^2 \dd t &\leq (1+x_0^2)\left(\underbrace{\int_0^T e^{Cte^{Ct}} \dd t}_{=: A(K,T)} - \int_0^T \frac{(Cte^{Ct})^{m+1}}{(m+1)!}\dd t\right)\\
        &= (1+x_0^2)\left(A(K,T) - \frac{1}{(m+1)!}\sum_{\ell=0}^{m+1}\binom{m+1}{\ell}\int_0^T (Ct)^{2(m+1)-\ell}\dd t\right)\\
        &= (1+x_0^2)\left(A(K,T) - \frac{1}{(m+1)!}\sum_{\ell=0}^{m+1}\binom{m+1}{\ell}\frac{T(CT)^{2(m+1)-\ell}}{2(m+1)-k+1}\right)\\
        &\leq (1+x_0^2)\left(A(K,T) - \frac{1}{(m+1)!}\sum_{\ell=0}^{m+1}\binom{m+1}{\ell}\frac{T(CT)^{2(m+1)-\ell}}{2(m+1)+1}\right)\\
        &= (1+x_0^2)\left(A(K,T) - \frac{1}{(m+1)!}\frac{T(CT)^{2(m+1)}}{2(m+1)+1}\left(1+\frac{1}{CT}\right)^{m+1}\right).
    \end{align*}
    By definition of $k^*$ it thus follows for all $(m,n) \in J_p$ that
    \begin{equation*}
        \int_0^T \left(\sum_{j=1}^p x_{k_j^*}(t)^2\right) \dd t \leq (1+x_0^2)\left(A(K,T) - \frac{1}{(m+1)!}\frac{T(CT)^{2(m+1)}}{2(m+1)+1}\left(1+\frac{1}{CT}\right)^{m+1}\right).
    \end{equation*}

    Combining the above results, we obtain
    \begin{align*}
        \hat E_\text{approx}^2 &\leq \E\left[\sum_{j=1}^p |\Psi_{k_j^*}-\widetilde{\Psi_j}|^2\right]\\ &\qquad\times \min_{(m,n) \in J_p}(1+x_0^2)\left(A(K,T) - \frac{1}{(m+1)!}\frac{T(CT)^{2(m+1)}}{2(m+1)+1}\left(1+\frac{1}{CT}\right)^{m+1}\right).
    \end{align*}
\end{proof}

\end{document}
